\SourceLeaf{44}entwickelt wird, lautet das allgemeine Glied $r^{\text{tum}} = \cos\left(2r+\frac{1-(-1)^n}{2}\right)\varphi$. Denn bei geradem $n$ vereinfacht sich diese Form offenbar zu $\cos 2r\varphi$, bei ungeradem geht sie in $\cos(2r+1)\varphi$ über.

In der kleinen Abhandlung selbst habe ich diese allgemeine Schreibweise nur in der abschließenden Formel (\S\,11) angewandt. Sie könnte zwar schon in den Untersuchungen, die zu dieser Formel führen, mit Nutzen angewandt werden. Ich hielt es jedoch für geboten, dort darauf zu verzichten, weil die nötigen Zusätze über die Umformung solcher Zeichen bei ihrem Eintritt in die Rechnungen selbst die Darstellung übermäßig verlängert und mich über die Grenzen der Abhandlung hinausgeführt hätten. Es schien mir zweckmäßiger, die Behandlung solcher Zeichen, wenn die von ihnen betroffenen Glieder in längere Rechnungen eingehen, und die für deren übersichtliche Durchführung erforderlichen Verfahren in einer eigenen Untersuchung zusammenzufassen.

1. Die Integration der vorgeschlagenen Formel scheint zunächst sehr leicht durchführbar. Denn die gegebene Größe $(1+n\cos\varphi)^\nu$ lässt sich nach dem binomischen Lehrsatz in eine Reihe nach aufeinanderfolgenden ganzzahligen Potenzen von $\cos\varphi$ umwandeln. Multipliziert man ihre einzelnen Glieder mit $d\varphi$ und integriert sie, so erhält man das Gesuchte. Das ist zwar richtig, doch führt dieser Weg zu sehr komplizierten Ausdrücken. Jeder Teil des gesuchten Integrals wird nämlich durch eine recht lange Reihe dargestellt. Daher empfiehlt sich ein anderes Verfahren. Es beruht darauf, dass sich nach den Regeln der analytischen Trigonometrie jede Potenz von $\cos\varphi$ mit positivem ganzzahligen Exponenten in eine sehr übersichtliche Reihe entwickeln lässt, deren Glieder nur die Kosinus von Vielfachen von $\varphi$ enthalten. Die gestellte Aufgabe wird also so gelöst, dass zuerst $(1+n\cos\varphi)^\nu$ durch eine Reihe nach Potenzen von $\cos\varphi$ dargestellt wird. Dann wird nach dem zuvor festgestellten allgemeinen Gesetz jedes ihrer Glieder in eine neue Reihe nach Kosinus von Vielfachen verwandelt. Schließlich fasst man aus allen diesen Reihen die Glieder zusammen, die den Kosinus desselben Winkels enthalten. Die daraus hervorgehende neue und gehörig geordnete Reihe lässt sich nach Multiplikation mit $d\varphi$ sofort integrieren.

2. Zuerst ist also das Gesetz zu untersuchen, nach dem eine Potenz eines gegebenen Kosinus durch die Kosinus von Vielfachen des Winkels dargestellt wird. Dieser Untersuchung liegt die elementare trigonometrische Formel\HistoricalNote{*)}{*)}{\textsc{Kaestner}, Astron. Abh., Erste Sammlung, I, 10.} zugrunde, nach der gilt
